\section{问题重述}

\subsection{问题背景}

${C_4}$ 烯烃是重要的化工原料，广泛应用于化工产品及医药生产。乙醇可作为制备 ${C_4}$ 烯烃的原料，其催化偶合反应性能受到催化剂组合与反应温度等条件的共同影响。题设催化剂组合主要由 Co 负载量、Co/SiO$_2$ 与 HAP 的装料比以及乙醇浓度条件构成。乙醇转化率反映原料的转化程度，${C_4}$ 烯烃选择性反映目标产物在生成物中的占比，二者共同决定 ${C_4}$ 烯烃收率。因此，研究各实验因素与反应性能之间的关系，并进一步确定适宜的催化剂组合和温度，对优化乙醇偶合制备 ${C_4}$ 烯烃的工艺条件具有实际意义。

题目提供两组实验数据。附件 1 给出了 A1～A14、B1～B7 共 21 种催化剂组合在不同温度下的乙醇转化率及各类产物选择性，其中 A1～A14 采用装料方式 I，B1～B7 采用装料方式 II；附件 2 给出了 350~$^{\circ}$C 下某一催化剂组合在同一次实验不同反应时间的测试结果。

\subsection{问题重述}

基于上述实验数据，需要围绕温度响应、因素影响、工艺优化与追加实验设计完成以下四项研究：

\noindent\textbf{问题一：温度响应与时间变化规律。}
针对附件 1 中每种催化剂组合，分别研究乙醇转化率和 ${C_4}$ 烯烃选择性随温度的变化关系；同时利用附件 2 的连续测试结果，分析 350~$^{\circ}$C 下乙醇转化率、${C_4}$ 烯烃选择性及其他主要产物选择性随反应时间的变化特征。

\noindent\textbf{问题二：催化剂组合及温度的影响。}
综合比较不同实验条件下的反应结果，探讨 Co 负载量、Co/SiO$_2$ 与 HAP 装料条件、乙醇浓度条件、装料方式及反应温度等因素与乙醇转化率、${C_4}$ 烯烃选择性之间的关系，分析不同因素对两项主要反应指标的影响大小及差异。

\noindent\textbf{问题三：${C_4}$ 烯烃收率优化。}
在相同实验条件下，以 ${C_4}$ 烯烃收率尽可能高为目标，选择适宜的催化剂组合与反应温度；并在温度严格低于 350~$^{\circ}$C 的约束下，进一步确定相应的较优催化剂组合和温度条件。

\noindent\textbf{问题四：新增实验方案设计。}
结合已有实验数据及前三问所得认识，在仅允许增加 5 次实验的条件下，合理设计新增实验的催化剂组合与反应温度，使有限实验能够为进一步认识因素影响和选择较优工艺条件提供有效信息，并对各实验方案给出详细理由。
